%% ma: L10-B02-0017
%% lop: 10
%% chuong: 1
%% bai: 2
%% dang: 10.2.9
%% loai: DS
%% muc: [TH, TH, TH, VD]
%% dap_an: "ĐĐSĐ"
%% nguon: "Ảnh giáo viên gửi 10/10/2026 (ThS. Nguyễn Hoàng Việt, Chương 2 Tập hợp), A-Đề 1, Phần 2 Câu 15"
%% kien_thuc: "Bài 2 (tập hợp và các phép toán trên tập hợp)"
%% trang_thai: chinh-thuc
\begin{ex}
Lớp $10B_1$ có $7$ học sinh giỏi Toán, $5$ học sinh giỏi Lý, $6$ học sinh giỏi Hóa, $2$ học sinh chỉ giỏi Toán và Lý, $3$ học sinh chỉ giỏi Toán và Hóa, $1$ học sinh chỉ giỏi cả Lý và Hóa, $1$ học sinh giỏi cả $3$ môn Toán, Lý, Hóa. Các mệnh đề sau đúng hay sai?
\choiceTF
{\True Số học sinh chỉ giỏi môn Toán là $1$ học sinh.}
{\True Số học sinh chỉ giỏi môn Lý là $1$ học sinh.}
{Số học sinh chỉ giỏi môn Hóa là $2$ học sinh.}
{\True Số học sinh giỏi ít nhất một môn (Toán, Lý, Hóa) là $10$ học sinh.}
\loigiai{Chỉ giỏi Toán: $7-2-3-1=1$. Chỉ giỏi Lý: $5-2-1-1=1$. Chỉ giỏi Hóa: $6-3-1-1=1\ne2$. Số học sinh giỏi ít nhất một môn: $1+1+1+2+3+1+1=10$.}
\end{ex}
